\subsection{I-43. Grassmann lines repair weight-three connectors}
\label{I-43}

\paragraph{Statement.}
Let $E$ be any field, and let $a_1,\ldots,a_n$ be distinct
nodes in $E$. Every nonzero homogeneous multilinear polynomial
$f$, having ROABP width at most three in some variable order
and satisfying
\[
 f(X+s\mathbf1)=f(X),\qquad 0\le n-2\deg f\le3,
 \tag{I.156}
\]
has nonzero pullback under the balanced two-pole map
\[
 R_i=t+\frac{z_1}{a_i-y_1}+\frac{z_2}{a_i-y_2}.
 \tag{I.157}
\]
The order need not be known, and the assertion holds in every
characteristic. Together with the arbitrary-variable cubic
theorem P-19, this proves the same conclusion for every
nonzero multilinear width-at-most-three polynomial on at most
eleven variables, without a homogeneity or translation promise.

For total degree at most $D$, the finite hitting set in both
assertions has at most
\[
 (n+1)^2\binom{D+3}{3}
 \tag{I.158}
\]
points over a working field of size at least
$\max\{2n+1,D+1\}$, with the extension-input convention of
I-40. The new conclusion is for the balanced rational map
with common center $t\mathbf1$. It does not assert the same
result for arbitrary center slopes or for the ordinary
unweighted centered-$\mathrm{SV}_2$ map.

\paragraph{Background input I-B10.}
We use the established description of a projective line in a
Grassmannian: its three-planes all contain a fixed two-plane
and lie in a fixed four-plane. The exact exterior-algebra
input is Ghorpade, Patil, and Pillai (2009), Lemma~3, printed
p.~4 of the author/arXiv version: two independent decomposable
$\ell$-vectors have decomposable sum exactly when their
associated $\ell$-planes intersect in dimension $\ell-1$.
That statement and proof are over an arbitrary field. The
highest-weight-zero and highest-weight-one inputs are the
degree-zero and degree-one cases of I-B09, with the conventions
and elementary Bruhat audit already recorded in I-40.

\paragraph{1. The precise connector problem.}
Work over an algebraic closure for the proof. The factor and
minimal-core reduction of I-41 applies unchanged: homogeneous
multilinear factors use disjoint variable sets, inherit the
width by specializing complementary factors, and each is
translation-invariant. Their nonnegative weights add to the
weight on the used coordinates. Thus a counterexample to
(I.156) would have an irreducible factor of weight at most
three. I-41 covers weights at most two. We may assume that
$f$ is irreducible, depends on all its variables, and has
weight $m=n-2D=3$.

Construct its actual minimal affine cores from the coefficient
support inclusions, as in I-40, and pad consistently to width
three. Every proper bond rank is at least two, and the first
and last are exactly two. Every strict internal interval has
generic rank two or three; a rank-one polynomial outer
factorization would factor the scalar into two nonconstant
polynomials, as in I-37. In particular the first and last
internal factors are singular. I-40 bounds every run of
rank-three factors by three.

It suffices first to preserve the rank-two contractions
between consecutive singular factors at the common one-pole
center
\[
 \alpha_i=t+\frac{z}{a_i-y}.
 \tag{I.159}
\]
Here the two singular factors themselves are included in the
connector. Consecutive singular factors with no intervening
full-rank factor, or with one, involve at most three distinct
coordinates; I-38's all-characteristic triple dominance
preserves their nonzero compound matrices. A run of length
three is saturated and is covered by I-42, part~4.

Consider therefore two full-rank cores $A_{j+1},A_{j+2}$,
between singular anchors $A_j,A_{j+3}$. Let $\delta_j$ be the
determinant weight of the left support $L_j$. The two increments
of three and the total budget nine give
\[
 \delta_j\in\{0,1,2,3\},\qquad
 \delta'_{j+2}=3-\delta_j.
 \tag{I.160}
\]
Reversing the witnessing variable order exchanges left and
right support spaces and transposes this connector. Thus it
is enough to handle $\delta_j=0$ and $\delta_j=1$. This is
order reversal, not reciprocal multilinear reversal, and
does not change the generator.

\paragraph{2. A single invariant anchor is enough.}
Suppose $\delta_j=0$. The argument of I-41 identifies the
left support with the actual invariant submodule
$\Gamma^2V\subset V\otimes V$, its preceding support with
the standard two-dimensional module, and its anchor with
\[
 C(u)=\begin{pmatrix}1&u&0\\0&1&u\end{pmatrix},
 \qquad u=X_j.
 \tag{I.161}
\]
Normalize the two full-rank constant cores to the identity by
constant basis changes. Torus covariance gives
\[
 M(x_1,x_2)=
 \begin{pmatrix}1&A&B\\0&1&C\\0&0&1\end{pmatrix},
 \tag{I.162}
\]
where $A,C$ are linear in the two internal variables and
$B$ is quadratic. Write the arbitrary opposite affine anchor
as $D(v)$, with rows $D_1(v),D_2(v),D_3(v)$, where
$v=X_{j+3}$. Its generic matrix rank is two. Its three rows
are linearly independent as coefficient tensors in
$E[v]^{r_{j+3}}$: a constant relation between them would be
a relation between the three actual suffix coefficient
functions at the preceding minimal bond. The constant gauges
above preserve this independence.

The wedge of the two rows of $C(u)M D(v)$ is the row of its
second compound. Its coefficient on $u^2$ is exactly
\[
 D_2(v)\wedge D_3(v).
 \tag{I.163}
\]
If this polynomial vector is nonzero, substitute (I.159) and
consider the pole at $y=a_j$. Every other coordinate is
regular there. The coefficient of the double pole is a
nonzero scalar multiple of
$z^2[D_2(v_0)\wedge D_3(v_0)]$, where
$v_0=t+z/(a_{j+3}-a_j)$. This vector is nonzero because
$v_0$ is transcendental over the coefficient field. Hence
this connector retains rank two.

It remains to treat $D_2\wedge D_3=0$. A primitive rank-one
factorization over the univariate polynomial ring gives
\[
 D_2(v)=p_2(v)q(v),\qquad D_3(v)=p_3(v)q(v).
\]
Both rows are nonzero and linearly independent over $E$.
They are affine. If the primitive vector $q$ had positive
degree, the degree bound would force both $p_2,p_3$ to be
constants, a contradiction. Thus $q$ is a nonzero constant
row, $p_2,p_3$ are independent affine polynomials, and
$D_1(v)\wedge q$ is nonzero because $D$ has generic rank two.
Direct multiplication now gives the complete exterior row
\[
 \bigl[p_2(v)+(C+u)p_3(v)\bigr]
       \bigl[D_1(v)\wedge q\bigr].
 \tag{I.164}
\]
At $y=a_j$ its first factor has a nonzero simple-pole
coefficient proportional to $z p_3(v_0)$, whereas the second
is regular and nonzero. The terms $A,B$ cancel from this
exterior expression, so no other internal dependence can
cancel that pole. Again the connector retains rank two.
This proves the case $\delta_j=0$, without assuming that the
opposite anchor is invariant or self-dual.

\paragraph{3. Determinant weight one forces two equal positive weights.}
Suppose now that $\delta_j=1$. The nonzero Pl\"ucker vector
of $L_j$ is fixed by the upper unipotent and has torus weight
one. Its $\mathrm{SL}_2$ orbit spans the standard
two-dimensional module, by the highest-weight-one input.
Its projective orbit is therefore a line, all of whose
points are decomposable three-vectors. I-B10 gives fixed
spaces
\[
 K\subset L_j\subset T,\qquad \dim K=2,\quad \dim T=4.
 \tag{I.165}
\]
They are the intersection and span of all planes in that
orbit line, so both are $\mathrm{SL}_2$-stable.

We give the needed small-width consequence explicitly; it
does not assume semisimplicity in positive characteristic.
The parity of tensor weights and $\delta_j=1$ imply that
$j$ is odd. Thus the two torus weights of the stable space
$K\subset V^{\otimes j}$ are $q,-q$ for an odd positive
integer $q$, by Weyl symmetry. Take all contractions of its
last physical factor. Their span $W_K$ is an
$\mathrm{SL}_2$-stable subspace of the preceding actual
support $L_{j-1}$, and hence has dimension at most three.

Assume $q>1$, and choose weight vectors $k_+,k_-$ of $K$.
Their coefficient vectors on the last factor can have only
the four distinct weights
\[
 q-1,\ q+1,\ -q-1,\ -q+1.
 \tag{I.166}
\]
Weyl symmetry pairs the first with the last and the second
with the third. Nonzero coefficient vectors of distinct
weights are independent. Since $\dim W_K\le3$, exactly one
of those pairs can occur. Up to exchanging $e_0,e_1$, this
means
\[
 k_+=p\otimes e_0,\qquad k_-=r\otimes e_1,
 \tag{I.167}
\]
where $p,r$ are independent, or the same display with the
two final basis vectors exchanged.

Apply the unipotent that sends $e_0$ to $e_0+s e_1$ in
(I.167), and write its action on $p$ as $\rho(s)p$.
Since $K$ is stable, the resulting tensor
$\rho(s)p\otimes e_0+s\rho(s)p\otimes e_1$ must remain in
the span of the two displayed tensors. Its $e_0$ coefficient
forces $\rho(s)p$ into $E(s)p$; its $e_1$ coefficient then
forces the same nonzero vector into $E(s)r$, a contradiction.
For the exchanged case use the opposite unipotent. This
argument works over the rational function field in $s$ and
in every characteristic. Hence $q=1$.

The determinant character of the stable space $K$ is zero.
The quotient line $L_j/K$ therefore has weight one. The
stable two-dimensional quotient $T/K$ has weights $1,-1$,
because a rational $\mathrm{SL}_2$-module has trivial
determinant character. Consequently the weights on $L_j$ are
\[
 (1,1,-1).
 \tag{I.168}
\]
No splitting of either short exact sequence in (I.165) was
used, and this deduction excludes Frobenius-twisted larger
weights directly.

\paragraph{4. A unique internal pole detects the remaining connector.}
Choose a weight basis as in (I.168) and again normalize the
two full-rank constant cores to the identity. Torus covariance
allows their linear matrices only in the common
negative-to-positive weight step. They thus have the form
\[
 N_\ell=v_\ell e_3^{\mathsf T},\qquad
 (v_\ell)_3=0,\qquad N_\ell N_k=0.
 \tag{I.169}
\]
Their product is
$M=I+x_1N_1+x_2N_2$. Its second compound is affine in
$x_1,x_2$: the quadratic exterior term of a rank-one update
vanishes. Multiplying by the compounds of the two arbitrary
affine anchors gives
\[
 C_2\bigl(A_j(u)M(x_1,x_2)A_{j+3}(v)\bigr)
   =P_0(u,v)+x_1P_1(u,v)+x_2P_2(u,v).
 \tag{I.170}
\]
Here the $P_\ell$ are polynomial matrices and the displayed
matrix is not identically zero, because the original
connector has generic rank two.

If some $P_\ell$ with $\ell=1,2$ is nonzero, take the pole
at the node belonging to $x_\ell$. Every other summand is
regular. The coefficient of its simple pole is a nonzero
scalar multiple of $zP_\ell(u_0,v_0)$, where
\[
 u_0=t+\frac{z}{a_j-a_{j+\ell}},\qquad
 v_0=t+\frac{z}{a_{j+3}-a_{j+\ell}}.
 \tag{I.171}
\]
The two coefficients of $z$ are distinct, so $(u_0,v_0)$
is an invertible linear change of $(t,z)$. A nonzero
polynomial matrix $P_\ell$ therefore stays nonzero. If both
$P_1,P_2$ vanish, the nonzero matrix $P_0(u,v)$ survives by
pairwise dominance. This proves preservation in the
$\delta_j=1$ case and, by order reversal, in every case of
(I.160).

\paragraph{5. Simultaneous ranks, endpoints, and the scalar output.}
Every univariate singular-factor pivot is nonzero after
(I.159), since each $\alpha_i$ is transcendental. The
preceding arguments preserve the nonzero rank-two compound
of every connector between consecutive singular factors.
All these assertions hold as rational-function identities
in the same independent parameters $t,z,y$. They therefore
hold simultaneously; no independent seed is assigned to an
individual connector.

For clarity, use the same pivot-cleared rank factorizations
as I-36: a rank-two anchor is $U_iV_i$, with both factors
of full respective rank. A connector of rank two says that
the $2\times2$ contraction between its adjacent anchors is
invertible. Products of these invertible contractions
preserve every interval containing anchors; intervals made
only of full-rank cores remain invertible. Thus all strict
internal interval ranks are preserved. This deduction uses
the common evaluated pivots, not preservation of only
constant boundary functionals.

Minimality gives singular first and last internal factors.
Pairwise dominance preserves the nonzero endpoint/first-anchor
and last-anchor/endpoint contractions. The preserved internal
rank-two identifications propagate their nonzero active
projections, exactly as in I-41. Hence all the hypotheses of
I-42, part~3, are now proved at the same one-pole center.
The maximum invertible run is three, which is inside its
all-characteristic bound four. An assumed annihilation of
(I.157) gives the reset Hessian with its three tangent kernels;
I-42 kills that Hessian and then every output path. This
contradicts $f\ne0$. The factor reduction in part~1 proves
the full assertion (I.156).

\paragraph{6. All multilinear inputs on at most eleven variables.}
I-41 already handles at most ten variables. Suppose a
homogeneous width-three annihilator on eleven variables
exists. P-19 and the known minimum balanced degree leave
only degrees $4,5,6,7$. For degree four, every nonconstant
coefficient of common translation has degree at most three.
If any were nonzero, the coefficient of the highest power
of the translation parameter would retain width at most
three by the coefficient-cut leading-limit argument of I-40.
It would contradict P-19 or the absence of smaller degrees.
Thus the quartic is translation-invariant, has weight
$11-8=3$, and is excluded by the theorem just proved.
Multilinear reversal excludes degree seven.

For degree five, the same highest-translation-coefficient
argument would produce a width-three annihilator of degree
at most four, all of which have now been excluded. The
quintic is therefore translation-invariant of weight one,
and I-40 excludes it. Reversal also excludes degree six.
Passing to a lowest nonzero homogeneous component preserves
the witnessing width and proves the assertion for arbitrary
multilinear inputs on at most eleven variables.

Finally apply the denominator-safe interpolation in I-40:
after multiplication by
$\prod_i(a_i-y_1)(a_i-y_2)$, each selector degree is at most
$n$, and the total degree in the three amplitudes is at most
$D$. The selector grid outside all nodes and the amplitude
Newton simplex give exactly (I.158), with the stated field
and extension scope.\hfill$\square$

\paragraph{Provenance and remaining scope.}
The Grassmann-line fact is an established input, not a new
claim about Grassmannians. The new use here is its interaction
with the actual adjacent coefficient support of dimension at
most three and the resulting pole calculation for constrained
ROABP connectors. All endpoint, interval, and simultaneous
substitution assertions above concern the same parameter
field; none follows merely from preserving each local space
against constant rank-one functionals.

At weight four, determinant weights at the ends of an
unsaturated run can both exceed one. Those connectors are
not covered by the two classifications above. The unrestricted
balanced two-pair question, the next higher-degree finite-
variable cases, and the polynomial-in-growing-width isolation
gate remain separate. No historical priority claim is made.
